\section{Training without the solutions}\label{sec:labelfree}

\Cref{sec:norm} changed the norm and kept the labels. This section keeps the energy norm and
removes the labels: the label-free transformer is the same network trained on the energy of
the same 1,024 instances \eqref{eq:training}, without their solutions. It is compared with
the energy-norm-supervised transformer, from which it differs in the labels and also in the
weighting of the load cases \eqref{eq:lkgrad} and in the learning rate
(\cref{sec:methods:training}), and with the usual practice, supervision on the displacement
error.

\subsection{Three dimensions}\label{sec:labelfree:three}

The label-free transformer was more accurate than the displacement-supervised transformer,
with or without the energy term, on all five metrics, in the seed mean and in each of the
three seeds (\cref{tab:3d}): its displacement error was \numThreeFreeDispAdvantage{} lower,
its relative energy gap \numThreeAncOverFreeGap{} times lower than that of the variant with
the energy term, the better of the two, a factor of \numThreeAncOverFreeGapRoot{} in the
energy-norm error, and its von Mises error \numThreeFreeVm{} against \numThreeLabVm{} and
\numThreeAncVm{}. Compared instance by instance within each seed, it had the lower relative
energy gap on \numPairThreeLabGap{} of the \numThreePairs{} instance-seed pairs against the
displacement-supervised transformer, the lower von Mises error on \numPairThreeLabVm{} and
the lower displacement error on \numPairThreeLabDisp{}. The two transformers differ in the
objective and also in the learning rate, which was fixed in advance and not tuned for either
(\cref{sec:methods:training}). The pre-registered condition on the accuracy of the
label-free transformer required its displacement error to be within 10\% of the
displacement-supervised transformer's with 1,024 labels, and its relative energy gap to be
at least 40\% lower at every number of labels; it was met with a wide margin, at
\numGTwoParity{} and with a smallest advantage of \numGTwoAdvMin{}. It was one of three
conditions of the comparison's gate, whose plausibility condition an amendment written after the
initial run assessed from 64 labels only; under its registered form the gate would have read
NO-GO. The claims here rest on the accuracy condition and on the measurements, not on the
gate's verdict (supplementary Sections~\SuppPrereg{} and~\SuppThree{}). The comparison fixed
no noise guard; read after the run with the guard of the two-dimensional comparison,
\numThreeGuardWithin{} of the \numThreeGuardTotal{} differences of seed means lie within it,
two of them, in relative energy gap and peak stress error, through the instance of
\cref{sec:norm:three}.

Against the energy-norm-supervised transformer, which has the labels and the norm, the
label-free transformer was the less accurate. Removing the labels raised the median
energy-norm error by a factor of \numNlThreeLabelE{} and the median von Mises error by
\numNlThreeLabelVm{} (\cref{fig:normlabels}), against factors of \numNlThreeNormE{} and
\numNlThreeNormVm{} for changing the norm. The pre-registration of the energy-norm run fixed
this comparison as a reading without a criterion: the energy-norm-supervised transformer's
median relative energy gap was \numKdFreeRel{} lower than the label-free transformer's,
beyond the guard of \numKdFreeTau{}, and lower on \numKdFreePairs{} of the \numThreePairs{}
instance-seed pairs; with three seeds, the Welch 95\% interval of this reading runs from
\numKdFreeWelchLow{} to \numKdFreeWelchHigh{}. On the medians it was also the more accurate
in displacement (\numKdFreeDispRel{}), von Mises error (\numKdFreeVmRel{}) and peak stress
error (\numKdFreePeakRel{}), beyond the guard, and within the guard in critical-region
recall, readings among \numKdSecondary{} exploratory comparisons (supplementary
Section~\SuppThree{}). On the held-out instances the difference in relative energy gap was
\numKdHoldFreeRel{}, beyond a guard of \numKdHoldFreeTau{}. Which of the labels, the
weighting of the load cases and the learning rate accounts for it was not tested.

\subsection{Two dimensions}\label{sec:labelfree:two}

In two dimensions the label-free transformer was more accurate than the
displacement-supervised transformer on every metric but the displacement error, in each of
the three seeds (\cref{tab:2d}): its relative energy gap was \numCmLabOverFreeGap{} times
lower, a factor of \numCmLabOverFreeGapRoot{} in the energy-norm error, and its von Mises
error \numCmFreeVm{} against \numCmLabVm{}, but its displacement error was
\numCmFreeDispExcessAbs{} higher, and higher in each seed. Instance by instance it had the
lower relative energy gap on \numPairCmLabGap{} of the \numCmPairs{} instance-seed pairs and
the lower von Mises error on \numPairCmLabVm{}, but the lower displacement error on only
\numPairCmLabDisp{}; on \numPairCmOpposite{} of the pairs the two errors ranked the two
networks in opposite orders (\cref{rem:rank}). The pre-registered hypothesis H1, that the
label-free transformer has the lower relative energy gap, was supported: its seed mean was
\numHOneRel{} lower, beyond the guard of \numHOneTau{}, with the same readings on the
medians and in the intervals as H2a and H2b. We have not investigated why the label-free
transformer was the more accurate in displacement in three dimensions and not in two, with
the same code.

H3, the energy-norm-supervised against the label-free transformer, was pre-registered as a
reading without a criterion, and showed no difference beyond the guard. The
energy-norm-supervised transformer's relative energy gap, \numCmKnormGap{}, was
\numHThreeRel{} lower, within the guard of \numHThreeTau{}, which its seeds made wide
(\numCmKnormGapSeeds{}, against \numCmFreeGapSeeds{}). In every seed it was the more
accurate of the two on every metric, by \numHThreeSeedRel{} in relative energy gap in seeds
0, 1 and 2, and resampling the instances, with the seeds kept, put its relative energy gap
\numHThreeResLow{} to \numHThreeResHigh{} below the label-free transformer's (95\%
interval); the guard, which treats the three seeds as the replicates, does not separate the
two networks. On the medians, removing the labels raised the energy-norm error by a factor
of \numNlCmLabelE{} and the von Mises error by \numNlCmLabelVm{}, against factors of
\numNlCmNormE{} and \numNlCmNormVm{} for changing the norm (\cref{fig:normlabels}). On the
held-out instances H3 again showed no difference beyond the guard, and the label-free
transformer's seed-mean relative energy gap was the lower, \numHoldTwoFreeGap{} against
\numHoldTwoKnormGap{} (\cref{sec:holdout}).

\subsection{Fewer labels}\label{sec:labelfree:budget}

\begin{figure}[tbp]
\centering
\figfull{generated/fig_label_efficiency.pdf}
\caption{Seed-mean relative displacement error (top) and relative energy gap (bottom) of the
supervised networks against the number of labelled training instances, in three dimensions
(left) and two (right). Horizontal line: the label-free transformer, trained on the same
1,024 instances without their solutions.}
\label{fig:labeleff}
\end{figure}

With fewer labels the supervised networks were trained for fewer steps as well
(\cref{sec:methods:training}), and their errors were in general larger (\cref{fig:labeleff};
the values are in supplementary Section~\SuppThree{}). The figure shows the
displacement-supervised transformer with the energy term in three dimensions only, and the
energy-norm-supervised transformer in two only, since in three it was trained with 1,024
labelled instances only; the graph network was trained with 64 and 1,024 labels only. In three dimensions the label-free
transformer's relative energy gap was lower than that of every network trained on the
displacement error at every number of labels, and its displacement error lower than both
displacement-supervised transformers';
the graph network reached a lower displacement error than the label-free transformer with
1,024 labels, not with 64. The displacement-supervised transformer's relative energy gap did
not fall monotonically: its seed mean rose from 256 to 1,024 labels, through the instance of
\cref{sec:norm:three}. In two dimensions the label-free transformer's relative energy gap
was lower than those of the displacement-supervised transformer and the graph network at
every number of labels; its advantage over the displacement-supervised transformer fell from
\numCmAdvMax{} with 16 labels to \numCmAdvMin{} with 1,024. In displacement, the two
networks supervised on it needed 1,024 labels to pass the label-free transformer; supervised
in the energy norm, the transformer came within the guard of the label-free one on every
metric from 256 labels on (supplementary Section~\SuppTwo{}). With 16, 64 and 256 labelled
instances, \numCmLabWorseSixteen{}, \numCmLabWorseSixtyFour{} and \numCmLabWorseTwoFiveSix{}
of the displacement-supervised transformer's \numCmPairs{} instance-seed predictions were
worse than the zero field, against \numCmKnormWorseSixteen{}, \numCmKnormWorseSixtyFour{}
and \numCmKnormWorseTwoFiveSix{} for the energy-norm-supervised transformer and none for the
label-free one.
